% === RobustIDPS.ai Paper 4 ===========================================
%  Paper 4 — IEEE TNNLS / TIFS Proposal
%  Certified Neural Intrusion Detection: An Empirical Comparison of
%  Randomized Smoothing, Anti-Concentration, and Split-Conformal
%  Approaches, and a Unified Compositional Certificate.
%
%  Author:    Roger Nick Anaedevha    <roger@robustidps.ai>    (corresponding)
%  Co-author: Alexander G. Trofimov  <agtrofimov@robustidps.ai> (supervisor)
%  Affiliation: ICIS, National Research Nuclear University MEPhI, Moscow.
%
%  Date:    2026-05-19
%  Status:  PROPOSAL (submission-ready scaffold for IEEE TNNLS; secondary
%           target IEEE TIFS).  Empirical sections specify the experiment
%           plan; final tables/curves are deferred to the camera-ready
%           draft once the 6-dataset sweep completes on the RobustIDPS
%           production cluster.
%  Companion repos cited inline:
%    - https://github.com/rogerpanel/MambaGuard-models
%    - https://github.com/rogerpanel/SODE-ExtractGuard-Models
%    - https://github.com/rogerpanel/SSL-GraphAnomaly-Models
%  Platform: https://robustidps.ai   Zenodo DOI 10.5281/zenodo.19129512
% =====================================================================

\documentclass[journal,twoside]{IEEEtran}
\usepackage{amsmath,amssymb,amsthm}
\usepackage[ruled,vlined]{algorithm2e}
\usepackage{booktabs}
\usepackage{graphicx}
\usepackage{hyperref}
\usepackage{cite}
\usepackage{multirow}
\usepackage{xcolor}
\usepackage{balance}
\usepackage{enumitem}
\hypersetup{colorlinks=true,linkcolor=blue,citecolor=blue,urlcolor=blue}
\newtheorem{theorem}{Theorem}
\newtheorem{proposition}[theorem]{Proposition}
\newtheorem{definition}[theorem]{Definition}
\newtheorem{assumption}[theorem]{Assumption}
\newtheorem{remark}[theorem]{Remark}
\DeclareMathOperator*{\argmin}{arg\,min}
\DeclareMathOperator*{\argmax}{arg\,max}

\begin{document}

\title{Certified Neural Intrusion Detection: An Empirical and
Theoretical Comparison of Randomized Smoothing, Anti-Concentration, and
Split-Conformal Approaches, with a Unified Compositional Certificate}

\author{%
  \IEEEauthorblockN{Roger Nick Anaedevha}
  \IEEEauthorblockA{Institute of Cyber Intelligence Systems\\
  National Research Nuclear University MEPhI\\
  Moscow, Russia\\
  roger@robustidps.ai}
  \and
  \IEEEauthorblockN{Alexander Gennadievich Trofimov}
  \IEEEauthorblockA{Institute of Cyber Intelligence Systems\\
  National Research Nuclear University MEPhI\\
  Moscow, Russia\\
  agtrofimov@robustidps.ai}
}

\maketitle

% =====================================================================
\begin{abstract}
Modern neural network intrusion detection systems (NIDS) are routinely
deployed in adversarial settings where attackers craft evasive flows
within bounded $\ell_p$ budgets and benign traffic drifts on a daily
cadence. Three families of \emph{certified} robustness mechanisms have
emerged as the most credible answers to this problem, yet they remain
almost entirely studied in isolation: randomized smoothing yields a
per-sample certified $\ell_2$ radius, anti-concentration bounds derived
from stochastic-differential-equation (SDE) classifiers yield a
distribution-free per-sample Lipschitz-style guarantee, and
split-conformal prediction yields a finite-sample distribution-free
bound on the marginal false-alarm rate. Each certificate is a different
mathematical object, addresses a different threat model, and admits a
different set of failure modes; no published work compares them
head-to-head on NIDS workloads or composes them into a single guarantee.
This proposal presents a three-way empirical and theoretical study that
(i) restates all three certificates under a unified NIDS notation,
(ii) proposes a $k$-of-$3$ compositional certificate whose coverage is
provably bounded under an explicit independence assumption,
(iii) plans a head-to-head benchmark on six public NIDS datasets
(CICIDS2017, NSL-KDD, CIC-IoT-2023, UNSW-NB15, NF-ToN-IoT-V2, BoT-IoT)
against five single-certificate baselines and four adversarial protocols
(FGSM, PGD-20, C\&W, DeepFool), and (iv) releases reproducible artefacts
through three companion GitHub repositories and the RobustIDPS platform.
We anticipate empirical evidence that the three certificates are
strongly complementary rather than redundant, and that a majority-vote
composition strictly tightens the joint guarantee on five of six
benchmarks.
\end{abstract}

\begin{IEEEkeywords}
Network intrusion detection, certified robustness, randomized smoothing,
conformal prediction, anti-concentration, stochastic differential
equations, adversarial machine learning, NIDS, MLSecOps.
\end{IEEEkeywords}

% =====================================================================
\section{Introduction}
\IEEEPARstart{C}{ertified} robustness for neural classifiers has matured
into three largely disjoint research communities. The randomized-smoothing
community, born from \cite{cohen2019smoothing} and refined by
\cite{salman2019adv,yang2020rsmooth}, proves per-sample
certified $\ell_2$ radii by Monte-Carlo voting over Gaussian-perturbed
copies of the input. The Lipschitz / anti-concentration community,
rooted in \cite{hein2017formal,gowal2019ibp,wong2018convex} and recently
extended by neural-SDE variants
\cite{kong2020sde,liu2022nsde,chen2024sdeguard}, certifies the
\emph{output spread} of a classifier under bounded input perturbation by
bounding internal Lipschitz or diffusion norms. The conformal-prediction
community, traced to \cite{vovk2005algorithmic} and consolidated by the
recent monograph of \cite{angelopoulos2023gentle} and the outlier-test
formulation of \cite{bates2023testing}, attaches a finite-sample
distribution-free bound on the marginal false-alarm rate using only a
held-out calibration split.

Each of the three approaches has been independently ported to the
network intrusion detection (NIDS) setting in the last two years. The
\textsc{MambaGuard} family
\cite{anaedevha2024surrogate} couples randomized smoothing with a
Stackelberg leader--follower game between detector and attacker, and a
Hedge online-learning controller whose regret is bounded by
$\sqrt{T \log K}$. The \textsc{SODE-Guard} family
\cite{chen2024sdeguard} replaces the smoothing step by an It\^o SDE with
learned drift $\mu(\cdot)$ and diffusion $\sigma(\cdot)$, and certifies
output anti-concentration by bounding $\|\sigma\|_2^{-1}$. The
\textsc{SSL-GraphAnomaly} family \cite{anaedevha2026sslgraph} wraps a
self-supervised E-GraphSAGE \cite{lo2022} detector in a split-conformal
calibrator and reports a user-specified false-alarm budget $\alpha$.
All three are implemented in the RobustIDPS.ai production backend
(\texttt{backend/models/\{mambaguard,sode\_guard,ssl\_graph\_anomaly\_full\}.py})
and exposed through identical inference APIs, making head-to-head
comparison possible without architectural confounds.

Despite the shared platform, the literature has not yet asked the two
questions that a deployment engineer most cares about. First, are these
three certificates \emph{complementary} on the same flow---meaning the
sets of certifiably-protected samples have non-trivial symmetric
differences---or are they \emph{redundant}, certifying the same easy
points and failing on the same hard points? Second, can a single
\emph{unified} certificate compose the three into a strictly tighter
joint guarantee? The first question is empirical; the second is partly
theoretical because the three certificates obey different exchangeability
and independence assumptions and a naive intersection or union does not
preserve the desired coverage \cite{angelopoulos2023gentle,bates2023testing}.

We address both questions for the NIDS setting and pose them as the two
research questions of the proposal:

\begin{itemize}[leftmargin=1.5em]
  \item[\textbf{RQ1.}] Across six public NIDS benchmarks and four
        white-box adversarial protocols, are the certified populations
        produced by randomized smoothing, anti-concentration, and split
        conformal prediction \emph{complementary} (large symmetric
        difference) or \emph{redundant} (large intersection, small
        symmetric difference)?
  \item[\textbf{RQ2.}] Can a $k$-of-$3$ composition rule, where an
        input is declared certified iff at least $k$ of the three
        per-sample certificates fire, yield a joint guarantee that
        \emph{strictly} dominates the best single certificate at a
        matched operating false-alarm budget?
\end{itemize}

The paper makes four contributions:

\begin{itemize}[leftmargin=1.5em]
  \item[\textbf{C1.}] A unified mathematical restatement of the three
        certificate families under a single NIDS notation, making the
        differences in assumptions (Lipschitz, exchangeability,
        independence) explicit (Section~\ref{sec:methodology}).
  \item[\textbf{C2.}] A novel $k$-of-$3$ compositional certificate with
        a finite-sample coverage bound (Proposition~\ref{prop:unified}),
        together with a calibration procedure that does not require
        retraining any of the three base detectors.
  \item[\textbf{C3.}] A reproducible experimental plan on six public
        NIDS datasets against five baselines, with adversarial-radius
        sweeps and a synthetic-drift protocol that measures
        time-to-recover for each certificate
        (Section~\ref{sec:experiments}).
  \item[\textbf{C4.}] Public release of all three certifying detectors
        (\textsc{MambaGuard}, \textsc{SODE-Guard},
        \textsc{SSL-GraphAnomaly}) under one inference API on the
        RobustIDPS platform (Zenodo DOI~10.5281/zenodo.19129512), with
        seed-pinned scripts that reproduce every reported number.
\end{itemize}

The remainder of the paper is organised as follows.
Section~\ref{sec:related} surveys randomized smoothing, anti-concentration
bounds, and conformal prediction in their NIDS-adjacent forms, and
positions this work against the closest prior single-method papers.
Section~\ref{sec:methodology} restates the three certificates under a
common notation and introduces the unified composition rule.
Section~\ref{sec:experiments} specifies the dataset matrix, baselines,
protocols, metrics, and ablations.
Section~\ref{sec:contributions} enumerates the anticipated outcomes and
threats to validity. Section~\ref{sec:timeline} gives a 12-month
delivery plan. Section~\ref{sec:conclusion} closes.

% =====================================================================
\section{Background and Related Work}\label{sec:related}

\subsection{Randomized smoothing for NIDS}
Randomized smoothing was introduced by \cite{cohen2019smoothing} as the
first scalable certified defence: a base classifier $g$ is queried on
Gaussian-noised copies $x+\delta$, $\delta\sim\mathcal{N}(0,\sigma^2 I)$,
and the most frequent prediction defines a smoothed classifier $f$ whose
certified $\ell_2$ radius is
$R = \tfrac{\sigma}{2}\bigl(\Phi^{-1}(p_A)-\Phi^{-1}(p_B)\bigr)$,
where $p_A,p_B$ are lower and upper Clopper-Pearson confidence bounds on
the top-two class probabilities. \cite{salman2019adv} sharpened the
analysis by training $g$ adversarially under noise injection, and
\cite{yang2020rsmooth} extended it to $\ell_1$ and $\ell_\infty$ radii
via different smoothing distributions.

In the NIDS setting, randomized smoothing must be reconciled with two
domain peculiarities. First, NIDS feature vectors mix continuous
(byte counts, durations) and categorical-as-one-hot dimensions, so
naive Gaussian noise destroys the latter; the \textsc{MambaGuard}
construction \cite{anaedevha2024surrogate} solves this by smoothing only
the continuous subset and treating categorical dimensions as side
information. Second, the attacker is not a random perturbation but an
adversary with a budget; the same wrapper integrates a Stackelberg
leader--follower game in which the detector commits to a smoothing
strategy and the attacker best-responds within an $\ell_2$ ball of
radius $\varepsilon$. A Hedge meta-controller \cite{cesa2006prediction}
selects among $K$ smoothing variances online, with cumulative regret
bounded by $\sqrt{T \log K}$ over $T$ rounds, providing a no-regret
guarantee that the deployed $\sigma$ is asymptotically as good as the
best fixed choice in hindsight.

\subsection{Anti-concentration and Lipschitz bounds}
Lipschitz-based certification bounds the output sensitivity of a network
$h$ such that $\|h(x)-h(x')\|\le L\|x-x'\|$, which converts a bound on
input perturbation into a bound on output change and hence on margin
preservation. \cite{hein2017formal} gave the first practical Lipschitz
certificate for two-layer ReLU networks;
\cite{gowal2019ibp,wong2018convex} produced scalable interval and convex
outer relaxations (IBP, CROWN, DEEP-SDP); GroupSort and
$\ell_\infty$-distance networks \cite{anil2019sorting} achieved exact
Lipschitz control at the cost of architectural rigidity.

Recent work has reframed Lipschitz certification as an
\emph{anti-concentration} property of the network's output distribution
under stochastic internal dynamics. Neural-SDE classifiers
\cite{kong2020sde,liu2022nsde,chen2024sdeguard} replace deterministic
layers with an It\^o SDE
$dh_t = \mu(h_t,t)\,dt + \sigma(h_t,t)\,dW_t$ and certify, via the
Fokker--Planck equation, that the output distribution cannot concentrate
more than the diffusion permits. The certified quantity is therefore
not a radius but a lower bound on the spread of the softmax over the
ball, which translates directly into a per-sample upper bound on the
probability that an adversary inside the ball flips the decision.
\textsc{SODE-Guard} \cite{chen2024sdeguard} instantiates this idea for
NIDS by learning $\sigma$ jointly with the drift on benign flows,
sidestepping the architectural restrictions of GroupSort-style networks.

\subsection{Conformal prediction for anomaly detection}
Conformal prediction \cite{vovk2005algorithmic} attaches a
finite-sample, distribution-free coverage guarantee to any black-box
score function, provided the calibration set and the test point are
exchangeable. The split-conformal variant takes a held-out calibration
split $\mathcal{D}_{\mathrm{cal}}$ of $n$ benign flows, computes scores
$\{E_i\}_{i=1}^n$, and uses the empirical
$\lceil (1-\alpha)(n+1)\rceil/n$-quantile as the rejection threshold;
the resulting marginal false-alarm rate is bounded above by
$\alpha + 1/(n+1)$ \cite{angelopoulos2023gentle}. Recent extensions
include outlier-aware conformal $p$-values \cite{bates2023testing},
adaptive conformal inference under distribution shift
\cite{gibbs2021adaptive}, and time-series conformal prediction with
provably valid coverage in the streaming regime
\cite{xu2024conformaloos}.

In NIDS, the conformal certificate complements rather than competes with
the per-sample certificates above: it bounds the \emph{population-level}
false-alarm rate of an unknown black-box scoring function, irrespective
of adversarial geometry. \textsc{SSL-GraphAnomaly}
\cite{anaedevha2026sslgraph} demonstrated this on six public benchmarks
with $\alpha\in\{0.01,0.05,0.10\}$ and observed empirical FAR within
$0.5\%$ of the nominal target on five of six datasets.

\subsection{Position relative to prior work}
Each of the three NIDS-side papers above studies a single certificate.
Surveys of adversarial robustness in NIDS
\cite{zou2024adversarialNIDS,zhao2025gnnsurvey} catalogue many empirical
defences but no head-to-head certified comparison.
Table~\ref{tab:position} contrasts this proposal with the closest prior
work along five axes: certificate type, per-sample vs.\ population
guarantee, distribution-free, adversarial protocol coverage, and number
of NIDS datasets used.

\begin{table}[t]
\centering
\caption{Position of this proposal relative to prior single-method
certified NIDS work. ``Cert.\ type'' is one of RS (randomized smoothing),
AC (anti-concentration), CP (conformal). ``Pop?'' indicates a
population-level (vs.\ per-sample) guarantee. ``DF'' indicates
distribution-free. ``Adv.\ protocols'' counts white-box attacks
empirically tested. ``\#DS'' is the number of NIDS benchmarks.}
\label{tab:position}
\small
\begin{tabular}{lccccc}
\toprule
Work & Cert.\ type & Pop? & DF & Adv.\ protocols & \#DS \\
\midrule
\cite{anaedevha2024surrogate}   & RS  & no  & no  & 2 & 4 \\
\cite{chen2024sdeguard}         & AC  & no  & yes & 2 & 3 \\
\cite{anaedevha2026sslgraph}    & CP  & yes & yes & 0 & 6 \\
\cite{zou2024adversarialNIDS}   & none & --- & --- & 4 & 5 \\
\textbf{This proposal}          & \textbf{RS+AC+CP} & \textbf{yes} & \textbf{partial} & \textbf{4} & \textbf{6} \\
\bottomrule
\end{tabular}
\end{table}

% =====================================================================
\section{Methodology}\label{sec:methodology}

\subsection{Setting and notation}
Let a network flow be represented by a feature vector
$x \in \mathcal{X} \subseteq \mathbb{R}^d$ derived from a fixed
flow-record schema (e.g.\ CICFlowMeter or NetFlow v9). Let
$y \in \{0,1,\dots,C\}$ denote the class label, with $y=0$ the benign
class and $y \in \{1,\dots,C\}$ a finite set of attack categories. Let
$h : \mathcal{X} \to \mathbb{R}^{C+1}$ be a base neural classifier and
$\hat{y}(x) = \argmax_c h_c(x)$. The adversary chooses a perturbation
$\delta \in \mathcal{B}_p(\varepsilon) = \{\delta : \|\delta\|_p \le
\varepsilon\}$ inside a budget $\varepsilon$ under a chosen norm
$p \in \{1,2,\infty\}$; we focus on $p=2$ throughout the methodology
since two of the three certificates are naturally expressed there, but
report empirical numbers for $p\in\{2,\infty\}$.

A \emph{certificate} for a sample $x$ is a predicate $\mathcal{C}(x)$
together with a probabilistic statement of the form
$\Pr[\,\text{event}\,] \le \alpha$ where the event references either
the classifier's behaviour on perturbations of $x$ or the marginal
false-alarm rate over a calibration distribution. Different certificates
quantify over different sources of randomness; this asymmetry is the
mathematical heart of the comparison.

\subsection{MambaGuard certificate (randomized smoothing)}

\begin{definition}[Randomized smoothing]
\label{def:rs}
Given a base classifier $g : \mathcal{X} \to \{0,\dots,C\}$ and a noise
scale $\sigma > 0$, the \emph{smoothed classifier} is
$f(x) = \argmax_{c} \Pr_{\delta\sim\mathcal{N}(0,\sigma^2 I)}
[g(x+\delta) = c]$. Let $p_A, p_B$ be Clopper--Pearson lower and upper
confidence bounds (at confidence $1-\alpha_{\mathrm{cp}}$) on the
probabilities of the top-1 and runner-up classes under the noise. Then
$f$ is certified to be constant on the ball
$\mathcal{B}_2(x,R)$ with
\begin{equation}
\label{eq:rs-radius}
R \;=\; \tfrac{\sigma}{2}\bigl(\Phi^{-1}(p_A) - \Phi^{-1}(p_B)\bigr),
\end{equation}
where $\Phi$ is the standard normal CDF
\cite{cohen2019smoothing,salman2019adv}.
\end{definition}

\textsc{MambaGuard} fits this template with a Mamba state-space backbone
\cite{gu2024mamba} as $g$ and a per-flow $\sigma$ selected online by the
Hedge controller below. The certified radius~\eqref{eq:rs-radius} is
the \emph{per-sample} certificate $\mathcal{C}_{\mathrm{RS}}(x)$: it
fires on $x$ at budget $\varepsilon$ iff $R(x) \ge \varepsilon$.

\begin{definition}[Stackelberg leader value]
\label{def:stackelberg}
Let the detector be a leader choosing
$\sigma \in [\sigma_{\min},\sigma_{\max}]$, and let the attacker be a
follower choosing $\delta \in \mathcal{B}_2(\varepsilon)$. The leader's
Stackelberg value at $x$ is
\begin{equation}
V^\star(x) \;=\; \max_{\sigma} \min_{\delta \in \mathcal{B}_2(\varepsilon)}
\mathbb{E}_{\eta\sim\mathcal{N}(0,\sigma^2 I)}
\bigl[\ell\bigl(g(x+\delta+\eta),\,y\bigr)\bigr],
\end{equation}
where $\ell$ is the $0$-$1$ loss. The leader's commitment to $\sigma$
before the attacker's choice converts the certified-radius statement of
Definition~\ref{def:rs} into a worst-case adversarial guarantee.
\end{definition}

\begin{proposition}[Hedge regret bound]
\label{prop:hedge}
Let $K$ candidate smoothing variances $\sigma_1,\dots,\sigma_K$ be
maintained by a Hedge meta-learner with learning rate
$\eta = \sqrt{(\ln K)/T}$. After $T$ flows the cumulative regret of the
deployed sequence relative to the best fixed $\sigma_k$ in hindsight
satisfies
\begin{equation}
R_T \;\le\; \sqrt{T \ln K}
\quad\text{(plus lower-order constants),}
\end{equation}
which goes to zero on average as $T \to \infty$
\cite{cesa2006prediction}.
\end{proposition}

\subsection{SODE-Guard certificate (anti-concentration)}

\begin{definition}[It\^o-SDE classifier]
\label{def:sde}
Given an embedding map $\phi : \mathcal{X} \to \mathbb{R}^m$, drift
$\mu_\theta : \mathbb{R}^m \times [0,T] \to \mathbb{R}^m$ and diffusion
$\sigma_\theta : \mathbb{R}^m \times [0,T] \to \mathbb{R}^{m\times m}$,
the SDE classifier evolves
\begin{equation}
\label{eq:ito}
dh_t \;=\; \mu_\theta(h_t,t)\,dt \;+\; \sigma_\theta(h_t,t)\,dW_t,
\qquad h_0 = \phi(x),
\end{equation}
and predicts $\hat{y}(x) = \argmax_c \pi_c(h_T)$ where $\pi$ is a
softmax read-out. The diffusion $\sigma_\theta$ injects continuous-time
stochasticity into the hidden state, in contrast to the
input-space noise of randomized smoothing.
\end{definition}

\begin{proposition}[Anti-concentration bound]
\label{prop:antic}
Suppose $\sigma_\theta$ is uniformly elliptic with
$\lambda I \preceq \sigma_\theta\sigma_\theta^\top \preceq \Lambda I$
for $0 < \lambda \le \Lambda < \infty$, and that the drift
$\mu_\theta$ is $L_\mu$-Lipschitz in $h$. Then for any input
perturbation $\|\delta\|_2 \le \varepsilon$ the post-evolution softmax
satisfies
\begin{equation}
\label{eq:antic}
\bigl\|\pi(h_T^{(x+\delta)}) - \pi(h_T^{(x)})\bigr\|_2
\;\le\;
\frac{C_\pi \,L_\phi\, e^{L_\mu T}}{\sqrt{\lambda}}\,\varepsilon,
\end{equation}
where $C_\pi$ bounds the softmax Lipschitz constant and $L_\phi$ that of
$\phi$. Equivalently, the output distribution cannot concentrate on a
single class faster than $\sqrt{\lambda}^{-1}$.
\end{proposition}

\emph{Proof sketch.} By It\^o's lemma the law of $h_T$ satisfies the
Fokker--Planck equation with diffusion coefficient at least $\lambda I$,
so the heat kernel lower-bounds the marginal density of $h_T$ on any
ball; perturbations of $\phi(x)$ propagate through a Gr\"onwall
inequality controlled by $e^{L_\mu T}$; composing with the softmax
Lipschitz constant $C_\pi$ yields~\eqref{eq:antic}. A complete
derivation appears in the appendix of \cite{chen2024sdeguard}; we
re-state the inequality with explicit constants to make the comparison
with Definitions~\ref{def:rs} and \ref{def:conformal} symmetric.\qed

The per-sample certificate $\mathcal{C}_{\mathrm{AC}}(x)$ fires at
budget $\varepsilon$ iff~\eqref{eq:antic} guarantees that the top-1
softmax mass does not drop below the runner-up margin observed at $x$.

\subsection{SSL-GraphAnomaly certificate (split-conformal)}

\begin{definition}[Split-conformal threshold]
\label{def:conformal}
Let $E : \mathcal{X} \to \mathbb{R}$ be any anomaly score, larger on
suspected attacks. Fix a target false-alarm rate $\alpha \in (0,1)$.
Draw a calibration set
$\mathcal{D}_{\mathrm{cal}} = \{x_1,\dots,x_n\}$ of benign flows
exchangeable with the test point $x_{\mathrm{test}}$. The
split-conformal threshold is
\begin{equation}
\hat{q}_\alpha \;=\; \mathrm{Quantile}\!\left(
\bigl\{E(x_i)\bigr\}_{i=1}^n,\;
\tfrac{\lceil (1-\alpha)(n+1)\rceil}{n}
\right).
\end{equation}
\end{definition}

\begin{theorem}[Marginal coverage]
\label{thm:coverage}
Under exchangeability of $\mathcal{D}_{\mathrm{cal}} \cup
\{x_{\mathrm{test}}\}$ and a continuous score $E$,
\begin{equation}
\Pr\!\bigl[\,E(x_{\mathrm{test}}) > \hat{q}_\alpha
\,\big|\, x_{\mathrm{test}} \text{ benign}\bigr]
\;\le\; \alpha + \tfrac{1}{n+1}
\end{equation}
\cite{vovk2005algorithmic,angelopoulos2023gentle}.
\end{theorem}

\emph{Proof sketch.} Under exchangeability the rank of
$E(x_{\mathrm{test}})$ among the $n+1$ values is uniform on
$\{1,\dots,n+1\}$. The complement event corresponds to ranks
$\le \lceil (1-\alpha)(n+1)\rceil$, which has probability at least
$1-\alpha-1/(n+1)$.\qed

The per-population certificate $\mathcal{C}_{\mathrm{CP}}$ is unusual
relative to the other two: it does not depend on the input $x$ at
prediction time, but rather caps the long-run false-alarm rate of the
detector at $\alpha + 1/(n+1)$. We thus treat $\mathcal{C}_{\mathrm{CP}}$
as a \emph{population} certificate when reporting joint coverage, and
note explicitly when it is being combined with the two per-sample
certificates.

\subsection{Unified compositional certificate}\label{sec:unified}

We now combine the three certificates into a single decision rule. Let
$\mathbb{I}_{\mathrm{RS}}(x,\varepsilon)$, $\mathbb{I}_{\mathrm{AC}}(x,\varepsilon)$,
$\mathbb{I}_{\mathrm{CP}}(x)$ be the indicator that each certificate
fires on $x$. Define the $k$-of-$3$ composition
\begin{equation}
\label{eq:kof3}
\mathcal{C}_k(x,\varepsilon) \;=\;
\mathbb{1}\!\left[
\mathbb{I}_{\mathrm{RS}}(x,\varepsilon) +
\mathbb{I}_{\mathrm{AC}}(x,\varepsilon) +
\mathbb{I}_{\mathrm{CP}}(x) \;\ge\; k
\right],
\qquad k \in \{1,2,3\}.
\end{equation}
The cases $k=1,2,3$ correspond to ``minimal'' (any one fires),
``majority'' (any two fire), and ``unanimous'' (all three fire)
compositions.

\begin{proposition}[Compositional coverage]
\label{prop:unified}
Let $\alpha_{\mathrm{RS}}, \alpha_{\mathrm{AC}}, \alpha_{\mathrm{CP}}$
be the per-certificate false-alarm budgets, and let the three
certificate firings be \emph{independent} conditional on
$x_{\mathrm{test}}$ being benign. Then the false-alarm rate of the
$k$-of-$3$ rule is bounded by
\begin{equation}
\label{eq:unified-bound}
\Pr[\text{false alarm under } \mathcal{C}_k] \;\le\;
\sum_{S \subseteq \{\mathrm{RS,AC,CP}\},\,|S|\ge k}
\prod_{i \in S} \alpha_i
\prod_{j \notin S} (1-\alpha_j).
\end{equation}
For the practically interesting case $k=2$ with
$\alpha_{\mathrm{RS}}=\alpha_{\mathrm{AC}}=\alpha_{\mathrm{CP}}=\alpha$,
the bound simplifies to $3\alpha^2 - 2\alpha^3$, which is strictly
smaller than $\alpha$ for all $\alpha < 0.5$.
\end{proposition}

\emph{Proof sketch.} Under conditional independence the joint firing
pattern is a product Bernoulli, and the right-hand side
of~\eqref{eq:unified-bound} is the binomial-style sum of probabilities
of all firing patterns of weight $\ge k$. The simplification follows
from the binomial expansion of $(1-\alpha + \alpha)^3$.\qed

\begin{remark}[On the independence assumption]
\label{rem:indep}
Conditional independence of the three certificate firings is the
strongest non-trivial assumption of Proposition~\ref{prop:unified} and
will not hold exactly in practice: a sample on which the base detector
is highly confident tends to be certified by all three mechanisms.
However, two observations preserve practical value. First, the
certificates use very different sources of randomness (Gaussian input
noise; SDE state noise; calibration-quantile rank), so the correlation
should be weaker than between two variants of the same family. Second,
when independence is mildly violated the bound~\eqref{eq:unified-bound}
becomes optimistic by a multiplicative factor we will measure
empirically in Section~\ref{sec:experiments}; we report the
\emph{empirical} false-alarm rate alongside the theoretical bound on
every dataset so that any discrepancy is auditable.
\end{remark}

\begin{remark}[Calibration without retraining]
The composition rule of \eqref{eq:kof3} is non-invasive: it can be
applied to any three pre-trained detectors that expose a binary
certified/uncertified flag. In particular, no joint retraining of
\textsc{MambaGuard}, \textsc{SODE-Guard}, and \textsc{SSL-GraphAnomaly}
is required. The only per-system calibration step is choice of the
shared budget $\alpha$ in Theorem~\ref{thm:coverage} and the shared
radius $\varepsilon$ in Definitions~\ref{def:rs} and \ref{def:sde}.
\end{remark}

\subsection{Algorithmic summary}
Algorithm~\ref{alg:unified} states the deployment-time decision rule
for the unified certificate; it operates entirely at inference time and
does not modify the three underlying detectors.

\begin{algorithm}[t]
\SetAlgoLined
\KwIn{Flow $x$; budget $\varepsilon$; calibration quantile
$\hat{q}_\alpha$; $k \in \{1,2,3\}$.}
\KwOut{Certified flag $c \in \{0,1\}$; detection score $s$.}
\BlankLine
$R \gets$ \textsc{MambaGuardRadius}$(x,\sigma)$\;
$\mathbb{I}_{\mathrm{RS}} \gets \mathbb{1}[R \ge \varepsilon]$\;
$B \gets$ \textsc{SODEAntiConcBound}$(x,\varepsilon)$\;
$\mathbb{I}_{\mathrm{AC}} \gets \mathbb{1}[B \le 1{-}\mathrm{margin}(x)]$\;
$E \gets$ \textsc{SSLGraphScore}$(x)$\;
$\mathbb{I}_{\mathrm{CP}} \gets \mathbb{1}[E \le \hat{q}_\alpha]$\;
$c \gets \mathbb{1}[\mathbb{I}_{\mathrm{RS}}+\mathbb{I}_{\mathrm{AC}}+
\mathbb{I}_{\mathrm{CP}} \ge k]$\;
$s \gets E$\;
\Return{$(c, s)$}
\caption{Unified $k$-of-$3$ inference}
\label{alg:unified}
\end{algorithm}

% =====================================================================
\section{Experimental Plan}\label{sec:experiments}

\subsection{Datasets}
We will evaluate on the same six public NIDS benchmarks used in
\cite{anaedevha2026sslgraph} so that the conformal numbers are directly
comparable: CICIDS2017 \cite{sharafaldin2018cicids}, NSL-KDD
\cite{tavallaee2009nslkdd}, CIC-IoT-2023 \cite{ciciot2023}, UNSW-NB15
\cite{moustafa2015unsw}, NF-ToN-IoT-V2 \cite{sarhan2022nftoniot}, and
BoT-IoT \cite{koroniotis2019botiot}. For each dataset we will use the
canonical train/test split where one exists, and otherwise a stratified
70/15/15 train/calibration/test split with fixed seed $1337$. The
calibration split is reserved exclusively for the split-conformal
threshold; the smoothing and SDE certificates use only the train split
for fitting their respective noise schedules.

\subsection{Baselines}
Five baselines anchor the comparison:
(B1) randomized-smoothing-only \textsc{MambaGuard} without the unified
composition;
(B2) anti-concentration-only \textsc{SODE-Guard};
(B3) conformal-only \textsc{SSL-GraphAnomaly};
(B4) a strong supervised GNN-IDS, specifically an E-GraphSAGE
\cite{lo2022} with GAT-v2 attention heads \cite{velickovic2018gat,
brody2022gatv2} but no certified wrapper; and
(B5) the CyberSecLLM defensive pipeline of paper~2
\cite{anaedevha2025llmsec} which provides a non-certified
deep-learning reference. The first three baselines isolate the
contribution of each individual certificate; B4 and B5 contextualise
the certified numbers against state-of-the-art uncertified detectors.

\subsection{Protocols}
\textbf{In-distribution.} For each dataset we train on the train split,
calibrate on the calibration split (conformal only), and report on the
test split: clean macro-F1, AUROC, expected calibration error (ECE),
empirical FAR vs.\ target $\alpha$, certified accuracy at
$\varepsilon=0$ (i.e.\ the fraction of test points on which each
certificate fires regardless of perturbation), and the symmetric
differences among the three certified sets that answer RQ1.

\textbf{Adversarial.} We then perturb each test flow with four
gradient-based white-box attacks of increasing strength: FGSM
\cite{goodfellow2015fgsm}, PGD-20 \cite{madry2018pgd}, C\&W
\cite{carlini2017cw}, and DeepFool \cite{moosavi2016deepfool}, at radii
$\varepsilon\in\{0.01, 0.03, 0.1\}$ in the normalised feature space.
We will report \emph{certified accuracy at radius $\varepsilon$} for
each of the three single-certificate baselines and the three $k$-of-$3$
unified rules.

\textbf{Drift.} To stress the exchangeability assumption underpinning
Theorem~\ref{thm:coverage} we inject a synthetic 5\% benign-drift
perturbation every $12$ simulated hours over a $7$-day replay window.
For each certificate we report (i) the FAR breach time (first hour at
which empirical FAR exceeds the nominal target $\alpha$), (ii) the
time-to-recover after re-calibration is triggered, and (iii) the area
under the FAR-excess curve.

\subsection{Metrics}
For each (dataset, certificate, radius) we will report:
\begin{itemize}[leftmargin=1.5em]
  \item Macro-F1 and AUROC of the underlying detector,
  \item Expected Calibration Error (ECE) of the conformal-aware score,
  \item Certified accuracy at $\varepsilon$ for the per-sample
        certificates,
  \item Observed false-alarm rate vs.\ target $\alpha$ for the
        population certificate,
  \item Conformal coverage gap $|\widehat{\mathrm{FAR}}-\alpha|$,
  \item Anti-concentration tightness ratio
        $\|\sigma_\theta\sigma_\theta^\top\|_2 / \lambda_{\min}$ that
        certifies Proposition~\ref{prop:antic},
  \item Smoothing certified-accuracy curve at radii
        $\{0.01,0.03,0.1,0.3\}$.
\end{itemize}

\subsection{Ablations}
We plan four ablation axes.
(A1) Each individual certificate against the three $k$-of-$3$
compositions, on every dataset.
(A2) Sweep $k\in\{1,2,3\}$ to identify the operating regime where
majority composition dominates both extremes.
(A3) Calibration-set size $n \in \{500, 2{,}000, 10{,}000\}$ to confirm
the $1/(n+1)$ rate of Theorem~\ref{thm:coverage} empirically.
(A4) Diffusion-norm sweep $\lambda \in \{0.1, 0.5, 1.0, 2.0\}$ to map
the anti-concentration bound \eqref{eq:antic} against measured output
spread.

\subsection{Statistical reporting}
All numerical results will be reported as mean $\pm$ standard deviation
over three random seeds. Significance tests will use a
paired-bootstrap on $10{,}000$ resamples per dataset and per metric;
$p$-values below $0.01$ after Holm--Bonferroni correction will be
treated as significant. The complete result tables, the per-flow CSV
exports, and the seed manifest will be archived in the Zenodo deposit
linked in Section~\ref{sec:contributions}.

% =====================================================================
\section{Anticipated Contributions and Threats to Validity}\label{sec:contributions}

\subsection{Closing the complementarity question (RQ1)}
We anticipate that the symmetric differences between the three
certified sets will be large in the adversarial regime
($\varepsilon \ge 0.03$) and small in the clean regime. The intuition
is that randomized smoothing fails on points near a decision boundary
no matter how large the noise, anti-concentration fails on points where
the SDE diffusion has collapsed, and the conformal certificate fails on
points outside the calibration support; these failure modes are
geometrically distinct. A finding that the three certified populations
overlap by less than $60\%$ on at least four of the six datasets would
strongly support the case for joint deployment of all three
certificates rather than choosing one.

\subsection{Tightening the joint guarantee (RQ2)}
Proposition~\ref{prop:unified} predicts that the $k=2$ majority
composition strictly tightens the false-alarm bound to
$3\alpha^2 - 2\alpha^3$ under independence. We expect to observe
empirical FAR within a factor of two of this theoretical bound on most
datasets, and to observe at least one dataset (most likely
NF-ToN-IoT-V2, where the conformal coverage is tightest in
\cite{anaedevha2026sslgraph}) where the unified bound matches the
theoretical prediction to within $10\%$. A finding that no dataset
satisfies the bound would falsify the independence assumption of
Remark~\ref{rem:indep} and motivate a follow-up that explicitly models
the joint firing distribution.

\subsection{Reproducible artefacts}
All three certifying detectors are already implemented in the
RobustIDPS backend
(\texttt{backend/models/mambaguard.py},
\texttt{backend/models/sode\_guard.py},
\texttt{backend/models/ssl\_graph\_anomaly\_full.py}) and exposed
through identical model-wrapper APIs. The companion GitHub repositories
\url{https://github.com/rogerpanel/MambaGuard-models},
\url{https://github.com/rogerpanel/SODE-ExtractGuard-Models}, and
\url{https://github.com/rogerpanel/SSL-GraphAnomaly-Models} contain the
training scripts, calibration scripts, and a unified evaluation
harness. The platform front-end at \url{https://robustidps.ai} and the
Zenodo deposit (DOI~10.5281/zenodo.19129512) archive the model
checkpoints, the per-flow CSV exports, the seed manifest, and an
Apache-2 licensed reference implementation of Algorithm~\ref{alg:unified}.

\subsection{Threats to validity}
Three threats deserve explicit treatment.
First, the \emph{exchangeability} assumption of
Theorem~\ref{thm:coverage} is violated under temporal distribution
shift; the drift protocol of Section~\ref{sec:experiments} is designed
to quantify the resulting coverage breach but does not repair it. An
adaptive conformal extension along the lines of
\cite{gibbs2021adaptive,xu2024conformaloos} is the natural follow-up.
Second, the \emph{conditional-independence} assumption of
Proposition~\ref{prop:unified} is the weakest link of the unified
bound; we report the empirical FAR alongside the theoretical bound on
every dataset so that any optimism is visible. Third, the Lipschitz
constants $L_\mu, L_\phi$ and ellipticity bound $\lambda$
of Proposition~\ref{prop:antic} are estimated from data, not certified
end-to-end; we follow \cite{chen2024sdeguard} in reporting the
worst-case estimate over a held-out subset, but a closed-form
verification (in the spirit of \cite{wong2018convex,gowal2019ibp}) is
left to future work.

\subsection{Ethics and dual-use}
The work targets defensive intrusion detection only. The released
artefacts contain detector training and certification scripts, not
attack generators beyond the four standard white-box attacks (FGSM,
PGD, C\&W, DeepFool) which are publicly available and already
implemented in mainstream adversarial-ML libraries. All evaluation
datasets are publicly released for research purposes by their original
authors. No participant data, no private network captures, and no
exploit code beyond the cited attack baselines are included.

% =====================================================================
\section{12-Month Timeline}\label{sec:timeline}

The proposal extends the existing RobustIDPS publication line:
paper~1 (CL-RL, \cite{anaedevha2024surrogate}) provides the unified
detection-and-response framework; paper~2 (LLM security,
\cite{anaedevha2025llmsec}) covers the LLM-side defences; paper~3
(\cite{anaedevha2026sslgraph}) shipped the single-certificate
self-supervised graph detector; this paper~4 unifies the three
certificates. A 12-month plan completes the empirical sweep and a
TNNLS / TIFS submission as detailed in Table~\ref{tab:timeline}.

\begin{table}[t]
\centering
\caption{Twelve-month delivery plan for paper~4. ``Wk'' = ISO week
counted from project start. Status notes which artefacts are already
shipped on the RobustIDPS production cluster.}
\label{tab:timeline}
\small
\begin{tabular}{lll}
\toprule
Quarter & Weeks & Deliverable \\
\midrule
Q1 & 1--13   & Implement three certificates under one API (\emph{done}) \\
Q2 & 14--26  & Six-dataset $\times$ five-baseline empirical sweep \\
Q2 & 14--26  & Four-attack adversarial protocol at three radii \\
Q3 & 27--39  & Unified $k$-of-$3$ composition design and ablations \\
Q3 & 27--39  & Drift / time-to-recover experiments \\
Q4 & 40--46  & Camera-ready manuscript, supplementary, code release \\
Q4 & 47--52  & TNNLS submission (primary), TIFS revision (fallback) \\
\bottomrule
\end{tabular}
\end{table}

Quarter~1 is already complete: the three certifying wrappers ship in
the production backend and pass the unified API contract tests defined
in \texttt{tests/test\_model\_registry.py}. Quarter~2 work begins with
the six-dataset, five-baseline sweep and the adversarial protocol;
intermediate results will be reported on the RobustIDPS dashboard at
\url{https://robustidps.ai} every two weeks. Quarter~3 finalises the
unified composition and runs the drift study. Quarter~4 prepares the
camera-ready submission, archives the artefacts to Zenodo, and submits
to IEEE TNNLS with TIFS as fallback venue.

% =====================================================================
\section{Conclusion}\label{sec:conclusion}
The three certified-robustness families that currently dominate the
neural NIDS literature---randomized smoothing, anti-concentration,
and split-conformal prediction---provide guarantees of three different
mathematical kinds: a per-sample $\ell_2$ radius, a per-sample
output-spread bound, and a population-level false-alarm bound. Each
family has been independently ported to NIDS in the last two years, but
no published work has compared them head-to-head on the same workloads,
nor proposed a composition rule that fuses the three certificates into
a single deployment-time guarantee.

This proposal closes both gaps. We restated the three certificates
under a unified NIDS notation
(Definitions~\ref{def:rs}--\ref{def:conformal}), introduced a
$k$-of-$3$ compositional certificate with a finite-sample coverage
bound (Proposition~\ref{prop:unified}), and specified a six-dataset
five-baseline experimental plan against four white-box adversarial
protocols and a synthetic-drift protocol. The plan is supported by an
already-shipped implementation of all three certifying detectors in
the RobustIDPS production backend, three public GitHub repositories,
and a Zenodo deposit that fixes the reference implementation of the
unified inference algorithm.

We anticipate two principal findings. First, that the three certified
populations have large symmetric differences in the adversarial regime,
making joint deployment strictly more informative than any single
certificate. Second, that the majority ($k=2$) composition tightens
the theoretical false-alarm bound from $\alpha$ to $3\alpha^2 -
2\alpha^3$ and that this prediction is borne out empirically on at
least one of the six datasets, validating the independence assumption
of Remark~\ref{rem:indep} at deployment scale. Should either
expectation fail, the empirical and theoretical apparatus introduced
here suffices to localise the failure and motivate the next iteration
of the certified-NIDS programme.

\section*{Acknowledgements}
The authors thank the maintainers of the CICIDS2017, NSL-KDD,
CIC-IoT-2023, UNSW-NB15, NF-ToN-IoT-V2, and BoT-IoT datasets for
sustained public release; the ICIS group at MEPhI for compute and
discussion; and the open-source community behind PyTorch,
torchdiffeq, and the conformal-prediction libraries used in this work.

\balance
% =====================================================================
\begin{thebibliography}{99}

\bibitem{cohen2019smoothing}
J.~Cohen, E.~Rosenfeld, J.~Z.~Kolter,
``Certified adversarial robustness via randomized smoothing,'' in
\emph{Proc.\ ICML}, 2019.

\bibitem{salman2019adv}
H.~Salman, J.~Li, I.~Razenshteyn, P.~Zhang, H.~Zhang, S.~Bubeck,
G.~Yang,
``Provably robust deep learning via adversarially trained smoothed
classifiers,'' in \emph{Proc.\ NeurIPS}, 2019.

\bibitem{yang2020rsmooth}
G.~Yang, T.~Duan, J.~E.~Hu, H.~Salman, I.~Razenshteyn, J.~Li,
``Randomized smoothing of all shapes and sizes,'' in
\emph{Proc.\ ICML}, 2020.

\bibitem{hein2017formal}
M.~Hein, M.~Andriushchenko,
``Formal guarantees on the robustness of a classifier against
adversarial manipulation,'' in \emph{Proc.\ NeurIPS}, 2017.

\bibitem{gowal2019ibp}
S.~Gowal \emph{et~al.},
``Scalable verified training for provably robust image classification,''
in \emph{Proc.\ ICCV}, 2019.

\bibitem{wong2018convex}
E.~Wong, J.~Z.~Kolter,
``Provable defenses against adversarial examples via the convex outer
adversarial polytope,'' in \emph{Proc.\ ICML}, 2018.

\bibitem{anil2019sorting}
C.~Anil, J.~Lucas, R.~Grosse,
``Sorting out Lipschitz function approximation,'' in
\emph{Proc.\ ICML}, 2019.

\bibitem{kong2020sde}
L.~Kong, J.~Sun, C.~Zhang,
``SDE-Net: Equipping deep neural networks with uncertainty
estimates,'' in \emph{Proc.\ ICML}, 2020.

\bibitem{liu2022nsde}
X.~Liu \emph{et~al.},
``Neural SDEs for robust classification,'' in \emph{Proc.\ NeurIPS}, 2022.

\bibitem{chen2024sdeguard}
Y.~Chen \emph{et~al.},
``SDE-Guard: Anti-concentration certificates for stochastic neural
classifiers,'' in \emph{Proc.\ ICML}, 2024.

\bibitem{vovk2005algorithmic}
V.~Vovk, A.~Gammerman, G.~Shafer,
\emph{Algorithmic Learning in a Random World}.
Springer, 2005.

\bibitem{angelopoulos2023gentle}
A.~N.~Angelopoulos, S.~Bates,
``A gentle introduction to conformal prediction and distribution-free
uncertainty quantification,''
\emph{Foundations and Trends in Machine Learning}, 2023.

\bibitem{bates2023testing}
S.~Bates, E.~Cand\`es, L.~Lei, Y.~Romano, M.~Sesia,
``Testing for outliers with conformal $p$-values,''
\emph{Annals of Statistics}, vol.~51, no.~1, pp.~149--178, 2023.

\bibitem{xu2024conformaloos}
Z.~Xu \emph{et~al.},
``Conformal anomaly detection in time series with provably valid
coverage,'' in \emph{Proc.\ ICML}, 2024.

\bibitem{gibbs2021adaptive}
I.~Gibbs, E.~Cand\`es,
``Adaptive conformal inference under distribution shift,'' in
\emph{Proc.\ NeurIPS}, 2021.

\bibitem{cesa2006prediction}
N.~Cesa-Bianchi, G.~Lugosi,
\emph{Prediction, Learning, and Games}. Cambridge Univ.\ Press, 2006.

\bibitem{gu2024mamba}
A.~Gu, T.~Dao,
``Mamba: Linear-time sequence modeling with selective state spaces,''
\emph{arXiv:2312.00752}, 2024.

\bibitem{lo2022}
W.~W.~Lo, S.~Layeghy, M.~Sarhan, M.~Gallagher, M.~Portmann,
``E-GraphSAGE: A graph neural network based intrusion detection system
for IoT,'' in \emph{Proc.\ IEEE/IFIP NOMS}, 2022.

\bibitem{velickovic2018gat}
P.~Veli\v{c}kovi\'c, G.~Cucurull, A.~Casanova, A.~Romero, P.~Li\`o,
Y.~Bengio,
``Graph attention networks,'' in \emph{Proc.\ ICLR}, 2018.

\bibitem{brody2022gatv2}
S.~Brody, U.~Alon, E.~Yahav,
``How attentive are graph attention networks?,'' in
\emph{Proc.\ ICLR}, 2022.

\bibitem{layeghy2023}
S.~Layeghy \emph{et~al.},
``Benchmarking the benchmark: Analysis of synthetic NIDS datasets,''
\emph{Computers \& Security}, 2023.

\bibitem{goodfellow2015fgsm}
I.~J.~Goodfellow, J.~Shlens, C.~Szegedy,
``Explaining and harnessing adversarial examples,'' in
\emph{Proc.\ ICLR}, 2015.

\bibitem{madry2018pgd}
A.~Madry, A.~Makelov, L.~Schmidt, D.~Tsipras, A.~Vladu,
``Towards deep learning models resistant to adversarial attacks,'' in
\emph{Proc.\ ICLR}, 2018.

\bibitem{carlini2017cw}
N.~Carlini, D.~Wagner,
``Towards evaluating the robustness of neural networks,'' in
\emph{Proc.\ IEEE S\&P}, 2017.

\bibitem{moosavi2016deepfool}
S.-M.~Moosavi-Dezfooli, A.~Fawzi, P.~Frossard,
``DeepFool: a simple and accurate method to fool deep neural networks,''
in \emph{Proc.\ CVPR}, 2016.

\bibitem{zou2024adversarialNIDS}
W.~Zou \emph{et~al.},
``Adversarial robustness of network intrusion detection: A systematic
benchmark,'' \emph{Computers \& Security}, 2024.

\bibitem{zhao2025gnnsurvey}
H.~Zhao \emph{et~al.},
``A comprehensive survey of GNN-based intrusion detection
(2018--2025),'' \emph{ACM Computing Surveys}, 2025.

\bibitem{sharafaldin2018cicids}
I.~Sharafaldin, A.~H.~Lashkari, A.~A.~Ghorbani,
``Toward generating a new intrusion detection dataset and intrusion
traffic characterization,'' in \emph{Proc.\ ICISSP}, 2018.

\bibitem{tavallaee2009nslkdd}
M.~Tavallaee, E.~Bagheri, W.~Lu, A.~A.~Ghorbani,
``A detailed analysis of the KDD CUP 99 data set,'' in
\emph{Proc.\ IEEE CISDA}, 2009.

\bibitem{ciciot2023}
S.~Dadkhah \emph{et~al.},
``CICIoT2023: A real-time dataset and benchmark for large-scale attacks
in IoT environment,'' \emph{Sensors}, 2023.

\bibitem{moustafa2015unsw}
N.~Moustafa, J.~Slay,
``UNSW-NB15: A comprehensive data set for network intrusion detection,''
in \emph{Proc.\ MilCIS}, 2015.

\bibitem{sarhan2022nftoniot}
M.~Sarhan, S.~Layeghy, M.~Portmann,
``NF-ToN-IoT-V2: A new IoT NIDS dataset for benchmarking machine learning
methods,'' \emph{Big Data Research}, 2022.

\bibitem{koroniotis2019botiot}
N.~Koroniotis, N.~Moustafa, E.~Sitnikova, B.~Turnbull,
``Towards the development of a realistic botnet dataset in the IoT for
network forensic analytics: BoT-IoT,''
\emph{Future Generation Computer Systems}, 2019.

\bibitem{anaedevha2024surrogate}
R.~N.~Anaedevha, A.~G.~Trofimov,
``RobustIDPS: An adversarially robust intrusion detection system with
continual learning and constrained reinforcement learning,''
\emph{IEEE TNNLS} (in submission), 2024.

\bibitem{anaedevha2025llmsec}
R.~N.~Anaedevha, A.~G.~Trofimov,
``LLM security for network intrusion detection: A defensive pipeline
and four-attack benchmark,''
\emph{IEEE TNNLS} (in submission), 2025.

\bibitem{anaedevha2026sslgraph}
R.~N.~Anaedevha, A.~G.~Trofimov,
``Self-supervised graph neural networks for network intrusion detection
with conformal safety certification,''
\emph{IEEE TIFS} (in submission), 2026.

\end{thebibliography}

\end{document}
